\documentclass{article}

\usepackage[dblblindworkshop,final]{neurips_2026}
\usepackage[utf8]{inputenc}
\usepackage[T1]{fontenc}
\usepackage{hyperref}
\usepackage{url}
\usepackage{booktabs}
\usepackage{amsfonts}
\usepackage{amsmath}
\usepackage{amssymb}
\usepackage{nicefrac}
\usepackage{microtype}
\usepackage[table]{xcolor}
\usepackage{graphicx}

\title{Lossless LRU Caching Is a Strong Baseline for Cache-Aware MoE Routing}
\workshoptitle{NeurIPS 2026 Workshop on Efficient Systems for Foundation Models}

\author{
  Pierre-Luc Bacon \\
  Mila, Quebec AI Institute; Universit\'e de Montr\'eal \\
  \texttt{PLACEHOLDER@umontreal.ca} \\[1em]
  Diego Calanzone \\
  Mila, Quebec AI Institute \\
  \texttt{diego.calanzone@mila.quebec} \\[1em]
  Riyasat Ohib \\
  Georgia Institute of Technology \\
  \texttt{PLACEHOLDER@gatech.edu}
}

\begin{document}
\maketitle

\begin{abstract}
Decoding a mixture-of-experts model is limited by memory movement, not arithmetic: when
experts overflow GPU memory, every routing decision that misses becomes a CPU-to-GPU copy.
Recent proposals learn to make this traffic temporally consistent -- hold one expert set for
many tokens, or reward a router directly for cache-friendly reuse -- and have mostly been
judged by proxies: switch rate, a simulated cache's hit ratio. We serve instead, on real
hardware with every copy paid on the decode path, and find that learned temporal structure
does not reliably translate into measured throughput: a controller of our own design and a
reward-fine-tuned router (CAR-MoE, Cache-Aware Routing for MoE, trained with GRPO against a
simulated LRU hit-ratio reward) both lose to simple, overlooked baselines on the metric that
matters -- physically counted expert transfers -- even where their own simulated instrument
reports gains. One measured constant, the per-expert transfer time $\tau$, predicts every
policy's decode cost within $\sim$2\%, and the real headroom is not where these methods look:
a clairvoyant eviction oracle moves $2.1\times$ fewer experts than LRU at zero quality cost,
locating the profitable learning problem in eviction, not in termination or reward shaping.
We extend this measurement to two additional backbones, Phi-tiny-MoE-instruct and
OLMoE-1B-7B, evaluating CAR-MoE, an option-critic-style controller, and MELINOE under an
independent, estimate-based protocol (compute plus a per-miss transfer penalty rather than a
real copy), and find a compatible story: our own routing-diversity analysis shows the
option-critic controller's near-perfect simulated hit ratio comes from collapsing the router
onto a handful of experts, at considerable cost to downstream task accuracy, while CAR-MoE
keeps routing close to load-balanced. Read together with the physical null result on OLMoE,
the message for this line of work is the same one switch-rate proxies already taught: a
metric not priced in moved bytes can improve while the served cost stands still.
\end{abstract}

\section{Introduction}

A sparse mixture-of-experts layer runs only a few experts per token, so most of its
parameters sit idle in GPU memory; OLMoE-1B-7B uses 1B of its 7B parameters per token
\citep{muennighoff2024olmoe}. When the model does not fit, the standard remedy keeps expert
weights in CPU memory and copies an expert to the GPU whenever a token's router selects one
that is not already resident, evicting whichever resident expert has gone unused longest
\citep{eliseev2023offloading}; the same design is entering mainline serving (vLLM pull
request \#37190, benchmarked on this very model). This LRU caching is \emph{lossless}: the
router always runs the experts it chose, so the model computes exactly what it would with
everything resident. Its price is traffic -- consecutive tokens select overlapping but not
identical experts, and decoding waits on the resulting stream of copies.

Two lines of recent work try to cut that stream by changing what the router does rather than
how eviction works. One draws on the option-critic framework \citep{bacon2017optioncritic}
and extends MoE routing temporally: hold an expert set as a macro-action until a learned
termination function decides to switch \citep{shen2026temporal}. The other treats
cache-friendliness as a reward and fine-tunes the router directly with policy gradients,
using a simulated LRU cache's hit ratio as the signal -- the approach we call CAR-MoE
(Cache-Aware Routing for MoE) in this paper, trained with GRPO
\citep{shao2024deepseekmath,yu2025dapo} rather than a bespoke auxiliary loss. Both have so
far been judged by proxies internal to their own instrumentation: switch rate, or the
simulated cache's own hit-ratio readout. Whether either buys any \emph{physical} serving
time over the lossless LRU baseline that actually moves weights has not been directly
measured until now.

We close that gap two ways. First, a serving harness keeps half of a frozen OLMoE-1B-7B's
expert weights in host memory only; every change to residency is a real, synchronous,
timed copy, and quality and traffic are measured in the same run. The verdict is that on
PCIe-4-class hardware, lossless per-token fetching beats every policy that holds a set
fixed, and the same harness, applied to a released CAR-MoE checkpoint, finds its simulated
hit-ratio gain does not reduce physically counted transfers at all. Second, we broaden the
comparison to two more backbones (Phi-tiny-MoE-instruct, OLMoE-1B-7B) and three objectives
(CAR-MoE, an option-critic-style controller, MELINOE \citep{raje2026melinoe}) under an
estimate-based protocol -- compute time plus a per-miss transfer penalty, calibrated against
each backbone's own base model -- paired with a routing-diversity analysis that asks not just
whether hit ratio improved, but how: by genuine temporal reuse of a load-balanced expert
population, or by routing collapse onto a fixed handful of experts. The two measurement
regimes disagree on whether CAR-MoE's simulated gain is real, and agree on where the actual
one-shot exploitable structure sits: eviction, not termination, and diversity-preserving
reuse, not collapse.

\section{Measurement Setup}
\label{sec:setup}

\paragraph{Physical harness (OLMoE-1B-7B, one backbone).} OLMoE-1B-7B, frozen. Eight of
sixteen MoE layers (the least dwell-sensitive under a per-layer probe) are restricted: their
64 experts live only in pinned host memory, served from a per-layer GPU pool of capacity
$C=32$. Decode is teacher-forced at batch 1 over sixteen held-out sequences; quality is
next-token NLL, throughput is measured with paired device events, and steady-state numbers
exclude each sequence's initial pool fill. A serving policy's cost per token is
\begin{equation}
  \text{ms/token} = c + \rho \cdot \tau,
  \label{eq:costmodel}
\end{equation}
where $c$ is compute time, $\rho$ the number of experts moved per token, and $\tau$ the
measured time to move one -- a hardware constant that spans $5\times$ across GPUs from PCIe~5
to PCIe~3. Calibrated by one measured $\tau$, Equation~\ref{eq:costmodel} reproduces measured
decode time within $\sim$2\% when copies are physically slowed $10\times$ and $50\times$, so
a single startup measurement of $\tau$ identifies the best policy at any tier.

\paragraph{Estimate-based cross-backbone protocol (Phi-tiny-MoE-instruct, OLMoE-1B-7B).} We
fine-tune both backbones with LoRA adapters (Phi: 16 experts, top-2; OLMoE: 64 experts,
top-8) on the Nemotron Post-Training Dataset v2 \citep{nvidia2025nemotron}, comparing base,
supervised fine-tuning (SFT) alone, CAR-MoE (a simulated fixed-capacity LRU working-set hit
ratio as the GRPO reward at one middle layer), an option-critic-style controller
\citep{shen2026temporal,bacon2017optioncritic}, and MELINOE. Unlike the physical harness
above, throughput here is \emph{not} a real copy: it is the base model's own measured
compute pass plus each objective's own working-set miss count times a benchmarked per-expert
disk-load constant, so an adapter that touches the expert projections directly is not
penalized twice for its own inference-time cost, but the copy itself is never actually
performed. We also compute the normalized routing entropy (ECI, $H(p)/\log(E)$, 1.0 =
load-balanced, near 0 = collapsed onto a few experts) to distinguish genuine temporal reuse
from routing collapse -- a distinction the physical harness's transfer count alone cannot
make. Every score is averaged over two independent seeds where both were available at
submission time.

\section{Results}

\subsection{The measured frontier holds the lossless baseline}

At the harness's native $\tau$, no policy that holds an expert set fixed beats exact LRU
without paying a quality cost of at least $+0.067$ nats: on PCIe-4-class hardware, temporal
extension does not pay. This is not a fixed verdict -- what carries across hardware tiers is
the transfer rate $\rho$, not the throughput, and Equation~\ref{eq:costmodel} converts any
$(\rho, c)$ pair to a throughput at any $\tau$. At $50\times$ slowed copies, the ordering
flips: low-traffic heuristics and the most switch-averse learned controller become the
fastest restricted policies. Selecting a policy is a lookup, not a fixed preference: measure
$\tau$ once, fix a quality tolerance, read the winner off Equation~\ref{eq:costmodel}.

\subsection{The Belady bound: the headroom is in eviction}

Replaying the recorded per-token expert demand through a clairvoyant oracle that always
evicts the expert needed furthest in the future \citep{belady1966} -- certified within 4\% of
the provable optimum -- yields exact serving (zero quality loss) at 4.6 transfers per token
against LRU's 9.6: a $2.1\times$ reduction for free. Holding a set fixed is therefore a
quality-priced approximation of what better eviction gives away; the profitable learning
problem is predicting near-future expert demand, not deciding when to stop holding a set.
Sixteen tokens of foresight recover 96\% of the oracle's advantage, and a causally-trained
scorer given only access history rediscovers a frequency heuristic -- the useful temporal
structure is already available losslessly, with a short horizon.

\subsection{Learned approaches under physical and estimated measurement}

\paragraph{The physical verdict on CAR-MoE.} The same harness, applied to a CAR-MoE
checkpoint (OLMoE-1B-7B fine-tuned with GRPO against a 16-slot simulated LRU hit-ratio
reward at one middle layer, with a prompt-conditioned variant selecting the budget via a
\texttt{[CACHE\_SIZE=c]} tag), finds no reduction in physically counted transfers relative to
the untuned base: 32,050 steady transfers against the base's 31,915, and 10,963 against
10,954 even at the exact layer and capacity CAR-MoE was trained on. The conditioning tag
makes traffic strictly worse. By its own instrument, training worked -- the simulated hit
rate at the trained layer improved -- but the counter serving real bytes did not move. This is
the same failure mode as the switch-rate proxy: a metric internal to the training-time
simulator, not priced in moved bytes, can improve while nothing served changes.

\paragraph{The estimate-based cross-backbone verdict.} Table~\ref{tab:mainresults} reports
our own protocol's downstream task deltas, working-set efficiency, and routing diversity
across both backbones. CAR-MoE is competitive with or better than the base model on most
tasks and improves hit ratio without collapsing the router (ECI stays within 0.01--0.02 of
the base on both backbones); combined with the supervised loss (Phi-tiny-MoE only, so far),
it raises hit ratio further while recovering the reward-only variant's estimated throughput
cost, which we read as evidence the reward combines readily with supervised learning rather
than requiring extended justification on its own. The option-critic controller instead
reaches a near-perfect simulated hit ratio by collapsing routing onto a handful of experts
(ECI falls from 0.89 to 0.28 on Phi-tiny-MoE), at a considerable cost to downstream accuracy
-- exactly the kind of instrument-internal gain the physical harness would need to confirm
independently, and which \S3.3's OLMoE re-measurement suggests should not be taken on faith.

\begin{table}[h]
\centering
\scriptsize
\caption{Estimate-based protocol, both backbones, at each backbone's trained working-set
capacity (Phi: $c=4$, top-2, 16 experts; OLMoE: $c=16$, top-8, 64 experts). Base row is
absolute; other rows are $\Delta$ from their own backbone's base. Avg~$\Delta$: mean of the
five task deltas. Misses: mean working-set misses per sequence (lower better). ECI:
normalized routing entropy (1.0 = load-balanced, 0 = collapsed). CAR-MoE rows highlighted;
OLMoE CAR-MoE~(+SFT) still running (--). Recall from \S3.3 that a physical re-measurement of
an OLMoE CAR-MoE checkpoint found no reduction in actual transferred bytes despite a
comparable simulated hit-ratio gain -- these numbers should be read as an internal,
proxy-based comparison across objectives, not as a claim about served throughput.}
\resizebox{\textwidth}{!}{%
\begin{tabular}{l cc cc cc cc cc cc | cc cc cc}
\toprule
 & \multicolumn{12}{c|}{lm-eval-harness score ($\Delta$ from base)} & \multicolumn{6}{c}{working-set efficiency} \\
 & \multicolumn{2}{c}{MMLU} & \multicolumn{2}{c}{MMMLU} & \multicolumn{2}{c}{GSM8K} & \multicolumn{2}{c}{HumanEval} & \multicolumn{2}{c}{MATH} & \multicolumn{2}{c|}{Avg $\Delta$} & \multicolumn{2}{c}{Hit ratio} & \multicolumn{2}{c}{Misses} & \multicolumn{2}{c}{ECI} \\
Objective & Phi & OL & Phi & OL & Phi & OL & Phi & OL & Phi & OL & Phi & OL & Phi & OL & Phi & OL & Phi & OL \\
\midrule
Base (untuned) & 0.62 & 0.56 & 0.34 & 0.38 & 0.62 & 0.68 & 0.03 & 0.28 & 0.09 & 0.06 & -- & -- & 0.57 & 0.69 & 649.8 & 1874.2 & 0.89 & 0.90 \\
\midrule
SFT & +0.00 & $-$0.02 & $-$0.02 & $-$0.01 & $-$0.02 & $-$0.09 & +0.01 & $-$0.13 & $-$0.00 & $-$0.02 & $-$0.01 & $-$0.05 & 0.56 & 0.68 & 742.1 & 1943.9 & 0.89 & 0.90 \\
Option-critic & $-$0.03 & $-$0.04 & $-$0.02 & $-$0.07 & $-$0.11 & $-$0.26 & $-$0.02 & $-$0.14 & $-$0.08 & $-$0.05 & $-$0.05 & $-$0.11 & 1.00 & 0.68 & 0.1 & 2012.7 & 0.28 & 0.89 \\
MELINOE & $-$0.00 & $-$0.07 & $-$0.01 & $-$0.06 & $-$0.12 & $-$0.31 & $-$0.02 & $-$0.20 & $-$0.07 & $-$0.05 & $-$0.05 & $-$0.14 & 0.88 & 0.76 & 237.7 & 1541.9 & 0.67 & 0.84 \\
\rowcolor{black!8} CAR-MoE & $-$0.01 & $-$0.00 & +0.00 & $-$0.01 & $-$0.05 & $-$0.01 & +0.00 & +0.05 & $-$0.01 & +0.00 & $-$0.01 & +0.01 & 0.62 & 0.76 & 435.9 & 1004.2 & 0.89 & 0.90 \\
\rowcolor{black!8} CAR-MoE (+SFT) & +0.00 & -- & +0.00 & -- & $-$0.03 & -- & +0.00 & -- & $-$0.01 & -- & $-$0.01 & -- & 0.78 & -- & 421.4 & -- & 0.71 & -- \\
\bottomrule
\end{tabular}%
}
\label{tab:mainresults}
\end{table}

\paragraph{Reading the two protocols together.} The physical harness and the estimate-based
protocol disagree on whether CAR-MoE's hit-ratio gain is real, because they measure different
things: the harness counts actual copies on one backbone at one trained configuration; our
protocol estimates a throughput number from compute time and a miss count that is never
physically moved, across two backbones and several objectives. Where they agree is more
important than where they differ: neither treats a simulated cache's own hit-ratio readout as
sufficient evidence of anything served, and the routing-diversity analysis available only in
our protocol shows that a high simulated hit ratio can come from two very different
mechanisms -- genuine, load-balanced temporal reuse (CAR-MoE) or router collapse onto a fixed
subset (option-critic) -- that a transfer count or a hit ratio alone cannot distinguish, but
that downstream task accuracy and a physical re-measurement both would.

\section{Related Work}

\citet{shen2026temporal} introduce the temporally extended mechanism judged by switch rate on
a resident model, with a deliberation cost against degenerate options
\citep{harb2018waiting,bacon2017optioncritic}; the missing number was the exchange rate from
switches to seconds, since their released implementation moves no weights and the memory
saving is arithmetic rather than physical. Systems that do move expert weights share the
design measured here -- experts in host memory, a fixed GPU cache, misses paid over the link --
and differ only in a hand-picked eviction rule, recency \citep{eliseev2023offloading} or
decayed frequency; none prices its rule against a measured hardware cost, and none learns
eviction. MELINOE \citep{raje2026melinoe} instead keeps the routing architecture and adds a
cache-simulation loss with a rank-matching regularizer against a frozen reference; we treat it
here as a third point in the same proxy-judged family as the option-critic controller and
CAR-MoE.

\section{Limitations and Conclusion}

The physical verdict is measured on one frozen 1B-active model, one text distribution, and
batch-1 teacher-forced decode, a regime batching dissolves; the cross-backbone protocol
substitutes an estimated cost model for a real copy, so its throughput numbers are a proxy in
exactly the sense the physical harness warns against, and should be read as an internal
comparison across objectives rather than a claim about served performance. Priced in real
bytes, cache-aware MoE routing is a hardware question: one measured constant and each
policy's transfer rate pick the winner at every tier, fast links keep the lossless baseline,
and proxy-aimed learning -- switch rate, a simulated cache's own hit ratio -- can buy nothing
the counter sees. Where CAR-MoE differs qualitatively from the option-critic controller is
not in whether its simulated gain is physically confirmed (it was not, on the one checkpoint
physically re-measured), but in \emph{how} it reaches that gain: without collapsing the
router's routing diversity, and combining readily with supervised fine-tuning. The eviction
oracle's factor-of-two headroom, not reward-shaped or architecturally hard-coded temporal
holding, remains where we think the profitable learning problem sits.

\bibliographystyle{plainnat}
\begin{thebibliography}{99}

\bibitem[Muennighoff et al.(2024)]{muennighoff2024olmoe}
Muennighoff, N., Soldaini, L., Groeneveld, D., Lo, K., Morrison, J., Min, S., Shi, W., Walsh, P., Tafjord, O., Lambert, N., et al.
\newblock OLMoE: Open Mixture-of-Experts Language Models.
\newblock \emph{arXiv preprint arXiv:2409.02060}, 2024.

\bibitem[Eliseev and Mazur(2023)]{eliseev2023offloading}
Eliseev, A. and Mazur, D.
\newblock Fast Inference of Mixture-of-Experts Language Models with Offloading.
\newblock \emph{arXiv preprint arXiv:2312.17238}, 2023.

\bibitem[Bacon et al.(2017)]{bacon2017optioncritic}
Bacon, P.-L., Harb, J., and Precup, D.
\newblock The Option-Critic Architecture.
\newblock \emph{Proceedings of the AAAI Conference on Artificial Intelligence}, 2017.

\bibitem[Shen and Henderson(2026)]{shen2026temporal}
Shen, Y. and Henderson, P.
\newblock Temporally Extended Mixture-of-Experts via Option-Critic.
\newblock \emph{arXiv preprint arXiv:2604.20156}, 2026.

\bibitem[Harb et al.(2018)]{harb2018waiting}
Harb, J., Bacon, P.-L., Klissarov, M., and Precup, D.
\newblock When Waiting is not an Option: Learning Options with a Deliberation Cost.
\newblock \emph{Proceedings of the AAAI Conference on Artificial Intelligence}, 2018.

\bibitem[Raje et al.(2026)]{raje2026melinoe}
Raje, A., Nayak, A., and Joshi, A.
\newblock MELINOE: Fine-Tuning Mixture-of-Experts Language Models for Cache-Friendly Offloaded Inference.
\newblock \emph{arXiv preprint arXiv:2602.11192}, 2026.

\bibitem[Shao et al.(2024)]{shao2024deepseekmath}
Shao, Z., Wang, P., Zhu, Q., Xu, R., Song, J., Bi, X., Zhang, H., Zhang, M., Li, Y.~K., Wu, Y., and Guo, D.
\newblock DeepSeekMath: Pushing the Limits of Mathematical Reasoning in Open Language Models.
\newblock \emph{arXiv preprint arXiv:2402.03300}, 2024.

\bibitem[Yu et al.(2025)]{yu2025dapo}
Yu, Q., Zhang, Z., Zhu, R., Yuan, Y., Zuo, X., Yue, Y., Fan, T., Liu, G., Liu, L., Liu, X., et al.
\newblock DAPO: An Open-Source LLM Reinforcement Learning System at Scale.
\newblock \emph{arXiv preprint arXiv:2503.14476}, 2025.

\bibitem[Belady(1966)]{belady1966}
Belady, L.~A.
\newblock A Study of Replacement Algorithms for a Virtual-Storage Computer.
\newblock \emph{IBM Systems Journal}, 5(2):78--101, 1966.

\bibitem[NVIDIA(2025)]{nvidia2025nemotron}
NVIDIA.
\newblock Nemotron Post-Training Dataset v2.
\newblock \emph{Hugging Face Dataset Card}, 2025.

\end{thebibliography}

\end{document}
